\chapter{The Stair That Was Not}

\section*{Told by}
Scholar Lyra, who never leaves the library and reads by candlelight. Her handwriting is too small to read without a lens. She says that is intentional. She also says the library has a ninth aisle, but no one has found it yet. She has been looking for forty years. The candles are getting shorter.

\section*{The Tale}

A student climbed a stair in the Collegium of Keys. He had counted the steps: one, two, three, four, five, six, seven, eight.

Then the ninth. There was no ninth. The stair continued anyway.

He should have stopped. Every architect in the Collegium could have told him: a stair that does not know how many steps it has is a stair that is not sure it exists. But he was young, and his thesis was due, and he had been awake for three nights. The body forgets how to be afraid when the mind is drowning.

He climbed until his legs ached. The walls changed from stone to vellum to glass --- the stone of the old foundations, the vellum of unbound manuscripts, the glass of a window that looks out onto a courtyard that has not existed since the fire of 1173. The steps became sentences. The banisters became arguments. The air smelled of ink and old mistakes.

At the top, a desk. No chair. No lamp. Just the desk and the dust. On it, a single page. It was his own thesis --- the one he had submitted three days ago, the one he had been so proud of, the one he had shown to his advisor who had nodded and said nothing.

It was covered in corrections. Not the advisor's hand. A hand he did not recognize. The letters were small --- smaller than his own, smaller than the lens could resolve. But he could read them. He could always read them. That was the worst part.

They were right. Every correction, every crossed-out line, every margin note that said \textit{``This is where you lied to yourself''} and \textit{``You knew this was wrong when you wrote it''} and \textit{``Your mother's death is not an argument, it is an excuse, and you know it.''}

He had been wrong. Not about the data --- the data was fine. He had been wrong about the question. He had asked the wrong thing. The stair had shown him what he should have asked instead.

He wept. Not because he was sad. Because the corrections were kind. Whoever had written them had taken the time to explain. Had drawn little diagrams in the margins. Had crossed out a paragraph and written, in letters no larger than a gnat's wing, \textit{``Try this instead. It will hurt less when you read it in the morning.''}

When he looked up, the desk was gone. The stair was gone. He was standing in the basement, holding a piece of chalk. The chalk was warm. It had been used recently.

He never told anyone what the corrections said. But his next thesis was perfect. Not clever --- perfect. The examiners had no questions. The advisors had no notes. The archivists bound it in leather and placed it on the highest shelf, where the dust falls upward and the light does not reach.

He never climbed a stair without counting twice. But he also never stopped writing. Every page he wrote after that night had a small mark in the corner --- a dot so tiny only he could see it. A reminder that somewhere, in a room that existed and did not exist, someone was reading. Someone was taking notes. Someone was being kind.

\section*{The Witch's Closing}

\textit{``A proof that cannot be improved is already dead. A scholar who cannot be wrong is a scholar who has stopped climbing. The stair judges --- but it also teaches. Accept its edits, or stay in the basement. The basement has no windows. The basement is where old arguments go to forget themselves.''}

\section*{What Lurks in the Shadow of the Tale}

\subsection*{The Invisible Stair}

A stair that appears only to those who have made a significant intellectual error --- not a small mistake, not a typo, but a wrong turn in the architecture of knowing. The kind of error that changes lives. The kind of error that, if uncorrected, will breed more errors, each one smaller and more poisonous than the last.

The stair leads to a room where the mistake is made manifest. The room is different for every climber. For some, it is a desk. For others, a mirror. For one architect in the Collegium's records --- a woman who designed a bridge that killed eleven people --- the room was a narrow ledge overlooking a river she had never seen. The water was cold. The ledge was narrow. She stood there for three hours before she spoke the words that let her leave.

To leave the room, the climber must admit the error aloud. Not to anyone --- the room has no witnesses except the climber and the stair. But the admission must be honest. The stair knows the difference between a confession and a performance.

Denial traps them. Not in the room --- the room is patient. It will wait. The trap is simpler: they cannot find the door. They can see it, a handspan away, but every step toward it moves them sideways. Some climbers have stayed for days. One stayed for a year. When he finally spoke the truth, he walked out and found that his wife had remarried and his children did not recognize his face. The stair did not apologize. The stair does not apologize. The stair corrects.

\paragraph{Tags / Mechanics:}
\begin{itemize}
\item [TRUTH] [DREAM] [THRESHOLD]
\item \textbf{Trigger:} A character fails a Lore, Investigation, or similar knowledge-based roll by 5 or more. The failure must be significant --- not a missed clue, but a genuine misunderstanding. The GM decides when the stair appears. It does not appear for every failure. Only for the ones that matter.
\item \textbf{To climb:} The character must succeed on a Spirit + Lore test (DV 6). On a success, they reach the room and see their error made manifest. On a failure, they climb but cannot reach the top --- the stair loops, the walls shift, and they emerge in a different part of the library (or dungeon, or forest) with no memory of what they were looking for. The stair does not forgive rushing.
\item \textbf{The Correction:} After admitting the error aloud (no roll required --- only honesty), the character gains +1 die to all future rolls of that skill. Not because they are smarter. Because they have seen the shape of their own wrongness. That shape is a kind of map.
\item \textbf{The Cost:} The stair does not give gifts for free. The character loses something tied to the original error: a relationship that depended on their certainty, a position they held because they were ``the one who knows,'' or a memory of the moment they first believed the wrong thing. The GM and player choose together. The loss should sting. It should not cripple.
\end{itemize}

\paragraph{Adventure Hook:}
A scholar's maps have been wrong for decades. She has hidden her error by avoiding the region entirely --- redirecting expeditions, burning letters, silencing anyone who asked too many questions. The region is a valley in the southern passes. The people there have been waiting for help that never came. The party must find the invisible stair, climb it, correct the map --- and then walk into the valley and face whatever the truth has done to the people who lived there. The scholar will not come with them. She is still in the basement. She has been there for three years. The chalk is almost gone.

\section*{Colophon}
Scriptorium of the Collegium, midnight. The scholar held up a lens. ``See the margin?'' she said. Her finger pointed to a blank space at the bottom of a page covered in notes. The space was not empty. There were letters there --- too small to read, even with the lens. ``That's where the stair begins. It begins everywhere. You just have to look closely enough to see the step.'' She blew out the candle. The room was dark. The stair was still there. It had always been there. It had been waiting for someone to stop counting and start climbing.